\chapter{Hyperparameter Configuration}
\label{appn:A}

This appendix provides a comprehensive reference of all hyperparameters used throughout the experiments presented in this report. Unless otherwise stated, all configurations were held constant across the three experimental conditions (model-free, imagination, and imagination with symlog) to ensure a fair comparison.

\section{Agent and Training Hyperparameters}

\Cref{tab:hyperparams_complete} lists every hyperparameter governing the agent architecture, optimisation procedure, and training loop.

\begin{table}[htbp]
\centering
\caption{Complete hyperparameter configuration used across all experiments.}
\label{tab:hyperparams_complete}
\begin{tabular}{@{}llp{6.2cm}@{}}
\toprule
\textbf{Hyperparameter} & \textbf{Value} & \textbf{Description} \\
\midrule
\multicolumn{3}{@{}l}{\textit{Optimisation}} \\
\addlinespace[2pt]
Learning rate           & $3 \times 10^{-4}$ & Adam optimiser, shared by all networks \\
Gradient clip norm      & 1.0                & Maximum gradient norm for clipping \\
Batch size              & 64                 & Mini-batch size sampled from replay buffer \\
\addlinespace[4pt]
\multicolumn{3}{@{}l}{\textit{Reinforcement Learning}} \\
\addlinespace[2pt]
Discount factor ($\gamma$) & 0.99           & Temporal discount applied to future rewards \\
Entropy coefficient     & 0.01               & Actor entropy bonus coefficient \\
\addlinespace[4pt]
\multicolumn{3}{@{}l}{\textit{Replay Buffer}} \\
\addlinespace[2pt]
Buffer size             & 100{,}000          & Maximum replay buffer capacity \\
\addlinespace[4pt]
\multicolumn{3}{@{}l}{\textit{Network Architecture}} \\
\addlinespace[2pt]
Hidden dimensions       & (256, 256)         & MLP hidden layer sizes (two layers) \\
\addlinespace[4pt]
\multicolumn{3}{@{}l}{\textit{World Model \& Imagination}} \\
\addlinespace[2pt]
Imagination horizon     & 5                  & Number of steps to imagine forward \\
Warmup steps            & 1{,}000            & Random action steps before imagination begins \\
WM train steps          & 2                  & World model gradient steps per training call \\
Imagination train ratio & 10                 & Imagination update performed every $N$ steps \\
\addlinespace[4pt]
\multicolumn{3}{@{}l}{\textit{Evaluation}} \\
\addlinespace[2pt]
Evaluation frequency    & 1{,}000            & Evaluate every $N$ environment steps \\
Evaluation episodes     & 10                 & Number of episodes per evaluation \\
\addlinespace[4pt]
\multicolumn{3}{@{}l}{\textit{System}} \\
\addlinespace[2pt]
Device                  & auto               & CUDA if available, otherwise CPU \\
\bottomrule
\end{tabular}
\end{table}

\section{Per-Environment Step Budgets}

Each environment was allocated a fixed total interaction budget, chosen to provide sufficient training time for convergence while remaining computationally tractable. \Cref{tab:step_budgets} summarises the step budgets used.

\begin{table}[htbp]
\centering
\caption{Total environment interaction steps allocated per environment.}
\label{tab:step_budgets}
\begin{tabular}{@{}lr@{}}
\toprule
\textbf{Environment} & \textbf{Total Steps} \\
\midrule
CartPole-v1    & 100{,}000 \\
LunarLander-v3 & 300{,}000 \\
\bottomrule
\end{tabular}
\end{table}

The step budget for \texttt{LunarLander-v3} is three times that of \texttt{CartPole-v1}, reflecting its higher-dimensional observation and action spaces and consequently slower convergence characteristics.

\section{Symlog Transform Implementation}

The symlog and symexp functions, adapted from \cite{hafner2023dreamerv3}, provide a bijective pair of transformations that compress the scale of target values while preserving their sign. \Cref{lst:symlog} presents the PyTorch implementations used in this project.

\begin{lstlisting}[
    language=Python,
    caption={Implementation of the symlog and symexp transforms and the associated loss function.},
    label={lst:symlog},
    basicstyle=\ttfamily\small,
    keywordstyle=\bfseries,
    commentstyle=\itshape,
    frame=single,
    numbers=left,
    numberstyle=\tiny,
    breaklines=true
]
def symlog(x):
    return torch.sign(x) * torch.log1p(torch.abs(x))

def symexp(x):
    return torch.sign(x) * (
        torch.exp(torch.clamp(torch.abs(x), max=20.0)) - 1
    )

def symlog_mse_loss(predictions, targets):
    return F.mse_loss(predictions, symlog(targets))
\end{lstlisting}

The \texttt{symlog} function compresses large magnitudes via $\operatorname{sign}(x) \cdot \ln(1 + |x|)$, ensuring that the world model's reward and value predictions operate on a numerically stable scale. The inverse \texttt{symexp} function recovers the original scale, with an internal clamp at $|x| \leq 20$ to prevent numerical overflow during exponentiation. The composite loss function \texttt{symlog\_mse\_loss} applies the symlog transform to the regression targets before computing the mean squared error, as described in \Cref{ch:method}.
